% !TEX root=../main.tex
\begin{abstract}

Fault diagnosis is critical in many domains, as faults may lead to safety threats or economic losses.
In the field of online service systems, operators rely on enormous monitoring data to detect and mitigate failures.
Quickly recognizing a small set of root cause indicators for the underlying fault can save much time for failure mitigation.
In this paper, we formulate the root cause analysis problem as a new causal inference task named \textit{intervention recognition}.
We proposed a novel unsupervised causal inference-based method named \textit{Causal Inference-based Root Cause Analysis} (CIRCA).
The core idea is a sufficient condition for a monitoring variable to be a root cause indicator, \textit{i.e.}, the change of probability distribution conditioned on the parents in the Causal Bayesian Network (CBN).
Towards the application in online service systems, CIRCA constructs a graph among monitoring metrics based on the knowledge of system architecture and a set of causal assumptions.
% \wenchi{for root cause analysis problem}
The simulation study illustrates the theoretical reliability of CIRCA.
The performance on a real-world dataset further shows that CIRCA can improve the recall of the top-1 recommendation by 25\% over the best baseline method.
% what the obstacles to deployment are
% However, a poor understanding of the system's normal status limits further improvement.

% Quickly recognizing a small set of indicators (namely \textit{root cause}) for the underlying fault can save much time for failure mitigation. \wenchi{is it necessary to state the relationship between the indicator and the root cause here? or only claim "recognizing a small set of relevant indicators for the ..." }
% In this paper, the root cause analysis problem is formulated as recognizing unexpected intervention to the system, and we find a sufficient condition for a monitoring variable to be one of the root cause indicators, \textit{i.e.}, the change of probability distribution conditioned on the parents in the Causal Bayesian Network.
% Hence, we proposed a novel unsupervised method named \textit{Structural Root Cause Analysis} (SRCA) to apply such a criterion in online service systems.
% \wenchi{In my view its more important to show novelty rather than mothod details in abstract}
% SRCA first constructs a graph among monitoring metrics based on the knowledge of system architecture and a set of causal assumptions \wenchi{for root cause analysis problem}.
% Taking the graph as a CBN, SRCA conducts hypothesis testing for each metric to check whether its distribution deviates from the expected one, instructed by the discovered criterion.
% The simulation study illustrates the theoretical reliability of SRCA, and the performance on a real-world dataset further shows SRCA can improves the recall of the top-1 recommendation by 51\% over the best baseline method.
 
 
\end{abstract}
